\documentclass[aspectratio=169]{beamer}
%--- Configuración del Tema y Colores ---
\usetheme{Madrid}
\usecolortheme{whale}
\setbeamertemplate{navigation symbols}{}
\setbeamertemplate{caption}[numbered]

%--- Paquetes ---
\usepackage[utf8]{inputenc}
\usepackage[spanish]{babel}
\usepackage{graphicx}
\usepackage{booktabs}
\usepackage{tikz}
\usepackage{fontawesome5}
\usepackage{amsmath}
\usetikzlibrary{shapes.geometric, arrows, positioning, calc, decorations.pathreplacing, patterns, fit}

%--- Metadatos del Documento ---
\title[PM2]{Programación Matemática 2}
\subtitle{Clase 12: Historia de la IA - El Surgimiento del Mecanismo de Atención}
\author[Luis Alfredo Alvarado]{Prof. Luis Alfredo Alvarado Rodríguez}
\institute[ECFM - USAC]{
    Escuela de Ciencias Físicas y Matemáticas \\
    Universidad de San Carlos de Guatemala
}
\date{Primer Semestre, 2026}

\begin{document}

%--- Diapositiva de Título ---
\begin{frame}
    \titlepage
\end{frame}

%--- Diapositiva: Tabla de Contenidos ---
\begin{frame}{Agenda de Hoy}
    \tableofcontents
\end{frame}

%===========================================================
% SECCIÓN 1: RECAPITULACIÓN
%===========================================================
\section{Recapitulación: El Problema de las Secuencias}

\begin{frame}{Lo que Vimos en la Clase 11}
    \centering
    \begin{tikzpicture}[scale=0.8]
        % Timeline
        \draw[thick, ->] (0,0) -- (12,0);
        
        \foreach \x/\year/\event in {
            1/1958/Perceptrón,
            3/1986/Backprop,
            5/1997/LSTM,
            7/2012/AlexNet,
            9/2014/Atención,
            11/2017/Transformer
        } {
            \draw (\x,0.15) -- (\x,-0.15);
            \node[below, font=\tiny] at (\x,-0.2) {\year};
            \node[above, font=\tiny, text width=1.3cm, align=center] at (\x,0.3) {\event};
        }
        
        % Highlight current topic
        \draw[red, thick, rounded corners] (8.3,-0.8) rectangle (11.7,1);
        \node[red, font=\scriptsize] at (10,-1.1) {Hoy: de aquí al Transformer};
    \end{tikzpicture}
    
    \vspace{0.5cm}
    \textbf{Resumen clase anterior:}
    \begin{itemize}
        \item CNNs revolucionaron visión por computadora (2012)
        \item RNNs procesan secuencias (texto, audio, series de tiempo)
        \item LSTMs solucionaron el problema de memoria a largo plazo
        \item Pero... aún había problemas fundamentales
    \end{itemize}
\end{frame}

\begin{frame}{RNNs: Procesamiento Secuencial}
    \centering
    \begin{tikzpicture}[scale=0.85]
        % RNN cells
        \foreach \i in {0,1,2,3,4} {
            \node[circle, draw, fill=yellow!30, minimum size=0.9cm] (h\i) at (\i*2.2,0) {$h_{\i}$};
            \node[circle, draw, fill=blue!20, minimum size=0.7cm] (x\i) at (\i*2.2,-1.5) {$x_{\i}$};
            \draw[->, thick] (x\i) -- (h\i);
        }
        
        % Recurrent connections
        \foreach \i in {0,1,2,3} {
            \pgfmathtruncatemacro{\j}{\i+1}
            \draw[->, thick, red] (h\i) -- (h\j);
        }
        
        % Labels
        \node[below, font=\scriptsize] at (0,-2) {El};
        \node[below, font=\scriptsize] at (2.2,-2) {gato};
        \node[below, font=\scriptsize] at (4.4,-2) {se};
        \node[below, font=\scriptsize] at (6.6,-2) {sentó};
        \node[below, font=\scriptsize] at (8.8,-2) {...};
        
        % Time arrow
        \draw[->, thick, gray] (0,-2.8) -- (9,-2.8) node[right, font=\scriptsize] {Tiempo};
    \end{tikzpicture}
    
    \vspace{0.4cm}
    \begin{alertblock}{Problema 1: Procesamiento Secuencial}
        Cada paso depende del anterior $\rightarrow$ \textbf{no se puede paralelizar}. Entrenar secuencias largas es muy lento.
    \end{alertblock}
\end{frame}

\begin{frame}{El Problema del Contexto Fijo}
    \centering
    \begin{tikzpicture}[scale=0.8]
        % Encoder
        \node[font=\scriptsize\bfseries] at (-1,1.5) {Encoder};
        \foreach \i in {0,1,2,3} {
            \node[circle, draw, fill=blue!20, minimum size=0.7cm] (e\i) at (\i*1.5,0) {};
        }
        \foreach \i in {0,1,2} {
            \pgfmathtruncatemacro{\j}{\i+1}
            \draw[->, thick] (e\i) -- (e\j);
        }
        
        % Context vector
        \node[rectangle, draw, fill=red!30, minimum width=0.8cm, minimum height=1.2cm] (ctx) at (6,0) {$\vec{c}$};
        \draw[->, thick] (e3) -- (ctx);
        
        % Decoder
        \node[font=\scriptsize\bfseries] at (11,1.5) {Decoder};
        \foreach \i in {0,1,2,3} {
            \node[circle, draw, fill=green!20, minimum size=0.7cm] (d\i) at (8+\i*1.5,0) {};
            \draw[->, thick, dashed] (ctx) -- (d\i);
        }
        \foreach \i in {0,1,2} {
            \pgfmathtruncatemacro{\j}{\i+1}
            \draw[->, thick] (d\i) -- (d\j);
        }
        
        % Labels
        \node[below, font=\tiny] at (0,-0.6) {The};
        \node[below, font=\tiny] at (1.5,-0.6) {cat};
        \node[below, font=\tiny] at (3,-0.6) {sat};
        \node[below, font=\tiny] at (4.5,-0.6) {down};
        
        \node[below, font=\tiny] at (8,-0.6) {El};
        \node[below, font=\tiny] at (9.5,-0.6) {gato};
        \node[below, font=\tiny] at (11,-0.6) {se};
        \node[below, font=\tiny] at (12.5,-0.6) {sentó};
        
        % Bottleneck annotation
        \draw[decorate, decoration={brace, amplitude=5pt}] (5.5,1.5) -- (6.5,1.5);
        \node[above, font=\tiny, red] at (6,1.7) {¡Cuello de botella!};
    \end{tikzpicture}
    
    \vspace{0.3cm}
    \begin{alertblock}{Problema 2: Vector de Contexto Fijo}
        Toda la oración de entrada debe comprimirse en UN solo vector $\vec{c}$. Para oraciones largas, se pierde información.
    \end{alertblock}
\end{frame}

\begin{frame}{El Problema de las Dependencias Largas}
    \centering
    \begin{tikzpicture}
        \node[font=\small, text width=12cm, align=left] at (0,0) {
            ``\textcolor{blue}{The cat}, which was sitting on the mat that my grandmother bought last summer when she visited the old furniture store downtown, \textcolor{blue}{was} very happy.''
        };
    \end{tikzpicture}
    
    \vspace{0.5cm}
    \begin{tikzpicture}[scale=0.8]
        % Words
        \node (w1) at (0,0) {The cat};
        \node (w2) at (3,0) {...20 palabras...};
        \node (w3) at (6,0) {was};
        
        % Connection
        \draw[->, thick, blue, bend left=40] (w1) to node[above, font=\tiny] {¿conexión?} (w3);
        
        % Gradient
        \draw[->, red, thick] (6,-0.8) -- (0,-0.8);
        \node[below, font=\tiny, red] at (3,-0.8) {Gradiente se desvanece};
    \end{tikzpicture}
    
    \vspace{0.3cm}
    \begin{alertblock}{Problema 3: Dependencias a Larga Distancia}
        Aunque LSTMs ayudan, las RNNs aún luchan para conectar información muy separada. El gradiente se ``diluye'' al retroceder muchos pasos.
    \end{alertblock}
\end{frame}

%===========================================================
% SECCIÓN 2: LA IDEA DE ATENCIÓN
%===========================================================
\section{La Idea Revolucionaria: Atención}

\begin{frame}{Inspiración: ¿Cómo Traducen los Humanos?}
    \begin{columns}
        \column{0.5\textwidth}
        \textbf{Enfoque RNN (incorrecto):}
        \begin{enumerate}
            \item Leer toda la oración en inglés
            \item ``Comprimir'' en la mente
            \item Escribir traducción al español
            \item Sin mirar atrás
        \end{enumerate}
        
        \vspace{0.3cm}
        \textbf{Enfoque Humano Real:}
        \begin{enumerate}
            \item Leer la oración
            \item Al traducir cada palabra...
            \item \textbf{Mirar de nuevo} las palabras relevantes
            \item Enfocarse donde es necesario
        \end{enumerate}
        
        \column{0.5\textwidth}
        \centering
        \begin{tikzpicture}[scale=0.7]
            % Source sentence
            \node[font=\scriptsize] (s1) at (0,3) {The};
            \node[font=\scriptsize] (s2) at (1,3) {black};
            \node[font=\scriptsize] (s3) at (2,3) {cat};
            \node[font=\scriptsize] (s4) at (3,3) {sat};
            
            % Target word being generated
            \node[font=\scriptsize, fill=yellow!30, rounded corners] (t1) at (1.5,0) {gato};
            
            % Attention weights (visual)
            \draw[thick, opacity=0.2] (s1) -- (t1);
            \draw[thick, opacity=0.3] (s2) -- (t1);
            \draw[thick, opacity=0.9, red] (s3) -- (t1);
            \draw[thick, opacity=0.1] (s4) -- (t1);
            
            % Attention label
            \node[font=\tiny, red] at (2.5,1.5) {Atención alta};
            
            \node[below, font=\scriptsize] at (1.5,-0.5) {``¿Qué palabras debo mirar?''};
        \end{tikzpicture}
    \end{columns}
\end{frame}

\begin{frame}{Bahdanau et al. 2014: El Paper Fundacional}
    \begin{block}{``Neural Machine Translation by Jointly Learning to Align and Translate''}
        Dzmitry Bahdanau, Kyunghyun Cho, Yoshua Bengio\\
        \textit{ICLR 2015 (arXiv 2014)}
    \end{block}
    
    \vspace{0.3cm}
    \textbf{La innovación clave:}
    \begin{itemize}
        \item En lugar de un vector de contexto \textbf{fijo}...
        \item Crear un vector de contexto \textbf{dinámico} para cada palabra de salida
        \item El modelo ``aprende a alinear'' entrada y salida
        \item Puede ``mirar atrás'' a cualquier parte de la entrada
    \end{itemize}
    
    \vspace{0.3cm}
    \begin{exampleblock}{Resultado}
        Mejora dramática en traducción automática, especialmente para oraciones largas.
    \end{exampleblock}
\end{frame}

\begin{frame}{Atención: La Intuición}
    \centering
    \begin{tikzpicture}[scale=0.85, transform shape]
        % Encoder hidden states
        \node[font=\scriptsize\bfseries] at (-0.5,2.5) {Encoder};
        \foreach \i/\word in {0/The, 1/cat, 2/is, 3/black} {
            \node[circle, draw, fill=blue!20, minimum size=0.8cm] (h\i) at (\i*2,2) {$h_{\i}$};
            \node[below, font=\tiny] at (\i*2,1.3) {\word};
        }
        
        % Current decoder state
        \node[circle, draw, fill=green!30, minimum size=1cm] (s) at (3,-1) {$s_t$};
        \node[below, font=\tiny] at (3,-1.8) {Decoder: generando ``negro''};
        
        % Attention weights
        \draw[thick, opacity=0.2] (h0) -- node[left, font=\tiny, opacity=1] {0.05} (s);
        \draw[thick, opacity=0.3] (h1) -- node[left, font=\tiny, opacity=1] {0.10} (s);
        \draw[thick, opacity=0.1] (h2) -- node[right, font=\tiny, opacity=1] {0.05} (s);
        \draw[thick, opacity=0.9, red] (h3) -- node[right, font=\tiny, opacity=1] {\textbf{0.80}} (s);
        
        % Context vector
        \node[rectangle, draw, fill=orange!30, minimum width=1.5cm, minimum height=0.6cm] (ctx) at (8,-1) {$c_t$};
        \draw[->, thick] (s) -- node[above, font=\tiny] {pesos} (ctx);
        
        % Annotation
        \node[font=\scriptsize, text width=3cm, align=center] at (8,-2.5) {Contexto dinámico\\(promedio ponderado)};
    \end{tikzpicture}
    
    \vspace{0.3cm}
    \[
    c_t = \sum_{i} \alpha_{ti} \cdot h_i \quad \text{donde} \quad \sum_i \alpha_{ti} = 1
    \]
\end{frame}

%===========================================================
% SECCIÓN 3: CÓMO FUNCIONA
%===========================================================
\section{¿Cómo Funciona la Atención?}

\begin{frame}{Los Tres Componentes: Query, Key, Value}
    \begin{columns}
        \column{0.5\textwidth}
        \textbf{Analogía: Búsqueda en una biblioteca}
        \begin{itemize}
            \item \textbf{Query (Q):} Lo que estás buscando\\
            \textit{``Quiero información sobre gatos''}
            \item \textbf{Key (K):} Etiquetas de los libros\\
            \textit{``Animales'', ``Historia'', ``Cocina''...}
            \item \textbf{Value (V):} El contenido de los libros\\
            \textit{El texto real que quieres leer}
        \end{itemize}
        
        \column{0.5\textwidth}
        \centering
        \begin{tikzpicture}[scale=0.8]
            % Query
            \node[rectangle, draw, fill=red!20, minimum width=1.5cm, minimum height=0.8cm] (q) at (0,0) {Query};
            
            % Keys
            \node[rectangle, draw, fill=blue!20, minimum width=1cm, minimum height=0.6cm] (k1) at (3,1.5) {$K_1$};
            \node[rectangle, draw, fill=blue!20, minimum width=1cm, minimum height=0.6cm] (k2) at (3,0.5) {$K_2$};
            \node[rectangle, draw, fill=blue!20, minimum width=1cm, minimum height=0.6cm] (k3) at (3,-0.5) {$K_3$};
            \node[rectangle, draw, fill=blue!20, minimum width=1cm, minimum height=0.6cm] (k4) at (3,-1.5) {$K_4$};
            
            % Values
            \node[rectangle, draw, fill=green!20, minimum width=1cm, minimum height=0.6cm] (v1) at (5,1.5) {$V_1$};
            \node[rectangle, draw, fill=green!20, minimum width=1cm, minimum height=0.6cm] (v2) at (5,0.5) {$V_2$};
            \node[rectangle, draw, fill=green!20, minimum width=1cm, minimum height=0.6cm] (v3) at (5,-0.5) {$V_3$};
            \node[rectangle, draw, fill=green!20, minimum width=1cm, minimum height=0.6cm] (v4) at (5,-1.5) {$V_4$};
            
            % Connections Q to K
            \draw[->, thick, opacity=0.3] (q) -- (k1);
            \draw[->, thick, opacity=0.8, red] (q) -- (k2);
            \draw[->, thick, opacity=0.2] (q) -- (k3);
            \draw[->, thick, opacity=0.4] (q) -- (k4);
            
            % K to V
            \foreach \i in {1,2,3,4} {
                \draw[-, dashed, gray] (k\i) -- (v\i);
            }
        \end{tikzpicture}
    \end{columns}
    
    \vspace{0.3cm}
    \textbf{Proceso:} Comparar Query con cada Key $\rightarrow$ obtener pesos $\rightarrow$ promediar Values
\end{frame}

\begin{frame}{El Cálculo de Atención (Simplificado)}
    \centering
    \begin{tikzpicture}[scale=0.75, transform shape]
        % Step 1: Scores
        \node[font=\small\bfseries] at (0,3) {Paso 1: Calcular scores};
        \node[rectangle, draw, fill=red!20] (q) at (0,2) {Q};
        \node at (1,2) {$\cdot$};
        \node[rectangle, draw, fill=blue!20] (k) at (2,2) {$K^T$};
        \node at (3,2) {$=$};
        \node[rectangle, draw, fill=purple!20] (scores) at (4.5,2) {scores};
        
        % Step 2: Softmax
        \node[font=\small\bfseries] at (8,3) {Paso 2: Normalizar};
        \node[rectangle, draw, fill=purple!20] (s2) at (7,2) {scores};
        \node at (8.2,2) {$\xrightarrow{\text{softmax}}$};
        \node[rectangle, draw, fill=orange!20] (weights) at (10,2) {pesos $\alpha$};
        
        % Step 3: Weighted sum
        \node[font=\small\bfseries] at (5,0.5) {Paso 3: Suma ponderada};
        \node[rectangle, draw, fill=orange!20] (w) at (3,-0.5) {$\alpha$};
        \node at (4,-0.5) {$\cdot$};
        \node[rectangle, draw, fill=green!20] (v) at (5,-0.5) {V};
        \node at (6,-0.5) {$=$};
        \node[rectangle, draw, fill=yellow!30] (out) at (7.5,-0.5) {Salida};
    \end{tikzpicture}
    
    \vspace{0.5cm}
    \[
    \text{Attention}(Q, K, V) = \text{softmax}\left(\frac{QK^T}{\sqrt{d_k}}\right) V
    \]
    
    \vspace{0.2cm}
    \begin{exampleblock}{Nota}
        El $\sqrt{d_k}$ es un factor de escala para estabilizar los gradientes. $d_k$ es la dimensión de los vectores.
    \end{exampleblock}
\end{frame}

\begin{frame}{Ejemplo Numérico Simple}
    \textbf{Supongamos vectores de 2 dimensiones:}
    
    \vspace{0.3cm}
    \begin{columns}
        \column{0.48\textwidth}
        \textbf{Entrada:}
        \begin{itemize}
            \item Query: $Q = [1, 0]$ (buscando)
            \item Keys: $K_1 = [1, 0]$, $K_2 = [0, 1]$
            \item Values: $V_1 = [10, 0]$, $V_2 = [0, 10]$
        \end{itemize}
        
        \column{0.48\textwidth}
        \textbf{Cálculo:}
        \begin{itemize}
            \item Score$_1 = Q \cdot K_1 = 1$
            \item Score$_2 = Q \cdot K_2 = 0$
            \item Softmax: $\alpha_1 \approx 0.73$, $\alpha_2 \approx 0.27$
        \end{itemize}
    \end{columns}
    
    \vspace{0.3cm}
    \textbf{Resultado:}
    \[
    \text{Output} = 0.73 \times [10, 0] + 0.27 \times [0, 10] = [7.3, 2.7]
    \]
    
    \vspace{0.2cm}
    \begin{block}{Interpretación}
        El Query ``se parecía más'' a $K_1$, así que la salida está dominada por $V_1$.
    \end{block}
\end{frame}

\begin{frame}{Visualización: Matriz de Atención}
    \centering
    \begin{tikzpicture}[scale=0.8]
        % Matrix
        \draw[step=0.8cm, gray, thin] (0,0) grid (4.8,4);
        
        % Row labels (target)
        \node[font=\tiny] at (-0.5,3.6) {El};
        \node[font=\tiny] at (-0.5,2.8) {gato};
        \node[font=\tiny] at (-0.5,2) {negro};
        \node[font=\tiny] at (-0.5,1.2) {duerme};
        \node[font=\tiny] at (-0.5,0.4) {.};
        
        % Column labels (source)
        \node[font=\tiny] at (0.4,4.4) {The};
        \node[font=\tiny] at (1.2,4.4) {black};
        \node[font=\tiny] at (2,4.4) {cat};
        \node[font=\tiny] at (2.8,4.4) {sleeps};
        \node[font=\tiny] at (3.6,4.4) {.};
        \node[font=\tiny] at (4.4,4.4) {EOS};
        
        % Attention weights (darker = higher)
        % Row 1: El -> The
        \fill[blue, opacity=0.8] (0,3.2) rectangle (0.8,4);
        \fill[blue, opacity=0.1] (0.8,3.2) rectangle (1.6,4);
        \fill[blue, opacity=0.1] (1.6,3.2) rectangle (2.4,4);
        
        % Row 2: gato -> cat
        \fill[blue, opacity=0.1] (0,2.4) rectangle (0.8,3.2);
        \fill[blue, opacity=0.2] (0.8,2.4) rectangle (1.6,3.2);
        \fill[blue, opacity=0.9] (1.6,2.4) rectangle (2.4,3.2);
        
        % Row 3: negro -> black
        \fill[blue, opacity=0.1] (0,1.6) rectangle (0.8,2.4);
        \fill[blue, opacity=0.9] (0.8,1.6) rectangle (1.6,2.4);
        \fill[blue, opacity=0.1] (1.6,1.6) rectangle (2.4,2.4);
        
        % Row 4: duerme -> sleeps
        \fill[blue, opacity=0.1] (1.6,0.8) rectangle (2.4,1.6);
        \fill[blue, opacity=0.85] (2.4,0.8) rectangle (3.2,1.6);
        
        % Row 5: . -> .
        \fill[blue, opacity=0.9] (3.2,0) rectangle (4,0.8);
        
        % Labels
        \node[font=\scriptsize, rotate=90] at (-1.5,2) {Salida (español)};
        \node[font=\scriptsize] at (2.4,5) {Entrada (inglés)};
    \end{tikzpicture}
    
    \vspace{0.3cm}
    Cada fila muestra ``a qué palabras de entrada atiende'' cada palabra de salida.
\end{frame}

%===========================================================
% SECCIÓN 4: EVOLUCIÓN
%===========================================================
\section{Evolución del Mecanismo de Atención}

\begin{frame}{Tipos de Atención}
    \begin{columns}[t]
        \column{0.32\textwidth}
        \begin{block}{Bahdanau (2014)}
            \textbf{Additive Attention}
            \[
            e_{ij} = v^T \tanh(W_1 h_i + W_2 s_j)
            \]
            \begin{itemize}
                \item Usa una red neuronal
                \item Más parámetros
                \item Primera propuesta
            \end{itemize}
        \end{block}
        
        \column{0.32\textwidth}
        \begin{block}{Luong (2015)}
            \textbf{Multiplicative Attention}
            \[
            e_{ij} = h_i^T W s_j
            \]
            \begin{itemize}
                \item Producto punto con matriz
                \item Más eficiente
                \item Varias variantes
            \end{itemize}
        \end{block}
        
        \column{0.32\textwidth}
        \begin{block}{Scaled Dot-Product}
            \textbf{(Transformer, 2017)}
            \[
            e_{ij} = \frac{q_i \cdot k_j}{\sqrt{d_k}}
            \]
            \begin{itemize}
                \item Solo producto punto
                \item Muy rápido
                \item Escalable
            \end{itemize}
        \end{block}
    \end{columns}
    
    \vspace{0.4cm}
    \begin{exampleblock}{Tendencia}
        La evolución fue hacia mecanismos más simples y eficientes computacionalmente.
    \end{exampleblock}
\end{frame}

\begin{frame}{De Atención como Complemento a Atención como Base}
    \centering
    \begin{tikzpicture}[scale=0.85, transform shape]
        % 2014-2016: Attention + RNN
        \node[font=\small\bfseries] at (0,3) {2014-2016};
        \draw[thick, rounded corners, fill=yellow!20] (-1.5,0) rectangle (1.5,2.5);
        \node at (0,2) {RNN/LSTM};
        \draw[thick, rounded corners, fill=red!20] (-1,0.5) rectangle (1,1.5);
        \node[font=\scriptsize] at (0,1) {+ Atención};
        \node[below, font=\tiny, text width=3cm, align=center] at (0,-0.3) {Atención como\\``complemento''};
        
        % Arrow
        \node at (3,1.25) {\faArrowRight};
        
        % 2017+: Only Attention
        \node[font=\small\bfseries] at (6,3) {2017+};
        \draw[thick, rounded corners, fill=red!30] (4.5,0) rectangle (7.5,2.5);
        \node at (6,1.25) {\textbf{Solo Atención}};
        \node[below, font=\tiny, text width=3cm, align=center] at (6,-0.3) {Atención es todo\\lo que necesitas};
        
        % Paper referencea
        \node[font=\scriptsize, text width=4cm, align=center] at (6,-1.2) {``Attention Is All You Need''\\Vaswani et al., 2017};
    \end{tikzpicture}
\end{frame}

\begin{frame}{El Paper que Cambió Todo (2017)}
    \begin{columns}
        \column{0.55\textwidth}
        \begin{block}{``Attention Is All You Need''}
            Vaswani, Shazeer, Parmar, Uszkoreit, Jones, Gomez, Kaiser, Polosukhin\\
            \textit{Google Brain / Google Research, 2017}
        \end{block}
        
        \vspace{0.3cm}
        \textbf{Propuesta radical:}
        \begin{itemize}
            \item Eliminar RNNs completamente
            \item Usar \textbf{solo} mecanismos de atención
            \item Arquitectura: \textbf{Transformer}
            \item Paralelización completa
        \end{itemize}
        
        \column{0.45\textwidth}
        \centering
        \begin{tikzpicture}[scale=0.7]
            % Citation count
            \draw[->] (0,0) -- (5,0) node[right, font=\tiny] {Año};
            \draw[->] (0,0) -- (0,4.5) node[above, font=\tiny] {Citas};
            
            \draw[thick, blue] (0.5,0.1) -- (1,0.3) -- (1.5,0.8) -- (2,1.5) -- (2.5,2.5) -- (3,3.5) -- (3.5,4);
            
            \node[font=\tiny] at (0.5,-0.4) {2017};
            \node[font=\tiny] at (2,-0.4) {2020};
            \node[font=\tiny] at (3.5,-0.4) {2024};
            
            \node[font=\tiny] at (4,4.2) {140,000+};
        \end{tikzpicture}
        
        \vspace{0.2cm}
        \scriptsize Uno de los papers más citados en la historia de CS
    \end{columns}
\end{frame}

\begin{frame}{¿Por Qué Fue Tan Importante?}
    \begin{columns}[t]
        \column{0.48\textwidth}
        \begin{block}{Problemas que Resolvió}
            \begin{itemize}
                \item[\faCheck] Paralelización total (no secuencial)
                \item[\faCheck] Dependencias de cualquier distancia
                \item[\faCheck] Contexto dinámico para cada posición
                \item[\faCheck] Entrenamiento mucho más rápido
            \end{itemize}
        \end{block}
        
        \column{0.48\textwidth}
        \begin{block}{Lo que Habilitó}
            \begin{itemize}
                \item BERT (2018)
                \item GPT (2018)
                \item GPT-2, GPT-3, GPT-4
                \item Claude, Gemini, Llama
                \item Toda la era de LLMs
            \end{itemize}
        \end{block}
    \end{columns}
    
    \vspace{0.4cm}
    \begin{alertblock}{Impacto}
        Sin el Transformer, no existirían ChatGPT ni ninguno de los LLMs modernos. El mecanismo de atención fue la clave.
    \end{alertblock}
\end{frame}

%===========================================================
% SECCIÓN 5: SELF-ATTENTION
%===========================================================
\section{Self-Attention: La Clave del Transformer}

\begin{frame}{Cross-Attention vs Self-Attention}
    \centering
    \begin{tikzpicture}[scale=0.8, transform shape]
        % Cross-Attention
        \begin{scope}[xshift=0cm]
            \node[font=\small\bfseries] at (1.5,3.5) {Cross-Attention};
            
            % Source
            \foreach \i in {0,1,2} {
                \node[circle, draw, fill=blue!20, minimum size=0.6cm, font=\tiny] (s\i) at (\i,2) {$h_\i$};
            }
            \node[font=\tiny] at (1,2.6) {Encoder};
            
            % Target
            \foreach \i in {0,1,2} {
                \node[circle, draw, fill=green!20, minimum size=0.6cm, font=\tiny] (t\i) at (\i,0) {$s_\i$};
            }
            \node[font=\tiny] at (1,-0.6) {Decoder};
            
            % Connections
            \foreach \i in {0,1,2} {
                \foreach \j in {0,1,2} {
                    \draw[->, gray, thin] (s\i) -- (t\j);
                }
            }
            
            \node[font=\tiny, text width=3cm, align=center] at (1,-1.5) {Atender de una secuencia a otra};
        \end{scope}
        
        % Self-Attention
        \begin{scope}[xshift=7cm]
            \node[font=\small\bfseries] at (1.5,3.5) {Self-Attention};
            
            % Single sequence
            \foreach \i in {0,1,2} {
                \node[circle, draw, fill=purple!20, minimum size=0.6cm, font=\tiny] (x\i) at (\i,1) {$x_\i$};
            }
            \node[font=\tiny] at (1,1.7) {Misma secuencia};
            
            % Self connections
            \draw[->, thick, red, bend left=30] (x0) to (x1);
            \draw[->, thick, red, bend left=30] (x0) to (x2);
            \draw[->, thick, red, bend left=30] (x1) to (x0);
            \draw[->, thick, red, bend left=30] (x1) to (x2);
            \draw[->, thick, red, bend left=30] (x2) to (x0);
            \draw[->, thick, red, bend left=30] (x2) to (x1);
            
            % Self loops
            \draw[->, thick, red] (x0) to [loop above] (x0);
            \draw[->, thick, red] (x1) to [loop above] (x1);
            \draw[->, thick, red] (x2) to [loop above] (x2);
            
            \node[font=\tiny, text width=3cm, align=center] at (1,-0.8) {Cada posición atiende a todas (incluyéndose)};
        \end{scope}
    \end{tikzpicture}
\end{frame}

\begin{frame}{Self-Attention: Cada Palabra ``Mira'' a Todas}
    \centering
    \begin{tikzpicture}[scale=0.9]
        % Sentence
        \node (w1) at (0,0) {El};
        \node (w2) at (1.5,0) {gato};
        \node (w3) at (3,0) {negro};
        \node (w4) at (4.5,0) {duerme};
        \node (w5) at (6.5,0) {tranquilo};
        
        % Attention from "negro"
        \node[above=0.3cm of w3, font=\scriptsize, red] {¿A quién atiendo?};
        
        \draw[->, thick, opacity=0.2] (w3) to [bend left=40] (w1);
        \draw[->, thick, opacity=0.9, red] (w3) to [bend left=40] (w2);
        \draw[->, thick, opacity=0.5] (w3) to [loop above] (w3);
        \draw[->, thick, opacity=0.3] (w3) to [bend right=40] (w4);
        \draw[->, thick, opacity=0.2] (w3) to [bend right=40] (w5);
        
        % Interpretation
        \node[below=1cm of w3, font=\scriptsize, text width=8cm, align=center] {
            ``negro'' atiende principalmente a ``gato'' porque es el sustantivo que modifica.\\
            El modelo \textbf{aprende} estas relaciones durante el entrenamiento.
        };
    \end{tikzpicture}
    
    \vspace{0.3cm}
    \begin{exampleblock}{Próximas Clases}
        En la \textbf{Clase 13} veremos la arquitectura completa del Transformer, y en la \textbf{Clase 14} profundizaremos en Self-Attention y Multi-Head Attention.
    \end{exampleblock}
\end{frame}

%===========================================================
% SECCIÓN 6: RESUMEN
%===========================================================
\section{Resumen y Ejercicios}

\begin{frame}{Línea de Tiempo: Del RNN al Transformer}
    \begin{table}[]
        \centering
        \renewcommand{\arraystretch}{1.2}
        \scriptsize
        \begin{tabular}{p{1.5cm} p{4cm} p{6.5cm}}
            \toprule
            \textbf{Año} & \textbf{Hito} & \textbf{Significado} \\
            \midrule
            1997 & LSTM & Memoria a largo plazo en RNNs \\
            2014 & Seq2Seq & Encoder-Decoder para traducción \\
            2014 & Bahdanau Attention & Primer mecanismo de atención neural \\
            2015 & Luong Attention & Variantes más eficientes \\
            2017 & Transformer & ``Attention Is All You Need'' \\
            2018 & BERT & Transformer bidireccional (encoder) \\
            2018 & GPT & Transformer autoregresivo (decoder) \\
            2020+ & GPT-3, etc. & Era de los LLMs masivos \\
            \bottomrule
        \end{tabular}
    \end{table}
\end{frame}

\begin{frame}{Conceptos Clave de Hoy}
    \begin{columns}[t]
        \column{0.48\textwidth}
        \begin{block}{Problemas de RNNs}
            \begin{itemize}
                \item Procesamiento secuencial (lento)
                \item Vector de contexto fijo
                \item Dependencias a larga distancia
            \end{itemize}
        \end{block}
        
        \begin{block}{Solución: Atención}
            \begin{itemize}
                \item Contexto dinámico
                \item Conexión directa a cualquier posición
                \item Pesos aprendidos
            \end{itemize}
        \end{block}
        
        \column{0.48\textwidth}
        \begin{block}{Fórmula Central}
            \[
            \text{Attn}(Q,K,V) = \text{softmax}\left(\frac{QK^T}{\sqrt{d}}\right)V
            \]
        \end{block}
        
        \begin{block}{Impacto}
            \begin{itemize}
                \item Habilitó el Transformer
                \item Base de todos los LLMs
                \item Revolucionó NLP y más allá
            \end{itemize}
        \end{block}
    \end{columns}
\end{frame}

\begin{frame}{Ejercicios Propuestos}
    \begin{enumerate}
        \setlength\itemsep{0.6em}
        \item \textbf{Cálculo manual:} Dados $Q=[1,0]$, $K_1=[0.5, 0.5]$, $K_2=[1, 0]$, $V_1=[10, 0]$, $V_2=[0, 10]$, calcular la salida de atención (sin el factor de escala).
        
        \item \textbf{Lectura:} Leer el abstract y la introducción del paper ``Attention Is All You Need''. Identificar los 3 problemas principales que mencionan sobre RNNs.
        
        \item \textbf{Visualización:} Usar la herramienta BertViz (\texttt{github.com/jessevig/bertviz}) para visualizar patrones de atención en un modelo pre-entrenado.
        
        \item \textbf{Reflexión:} ¿Por qué crees que se llama ``Self''-Attention? ¿Qué diferencia hay con la atención de Bahdanau?
        
        \item \textbf{(Bonus)} Implementar atención básica en NumPy: función que reciba Q, K, V y devuelva el resultado.
    \end{enumerate}
\end{frame}

\begin{frame}{Recursos Adicionales}
    \textbf{Papers fundamentales:}
    \begin{itemize}
        \item Bahdanau et al., ``Neural Machine Translation by Jointly Learning to Align and Translate'' (2014)
        \item Vaswani et al., ``Attention Is All You Need'' (2017)
    \end{itemize}
    
    \vspace{0.3cm}
    \textbf{Visualizaciones interactivas:}
    \begin{itemize}
        \item Jay Alammar: ``The Illustrated Transformer''
        \item 3Blue1Brown: ``Attention in Transformers''
    \end{itemize}
    
    \vspace{0.3cm}
    \textbf{Próximas clases:}
    \begin{itemize}
        \item Clase 13: Arquitectura Transformer completa
        \item Clase 14: Self-Attention en detalle matemático
    \end{itemize}
\end{frame}

%--- Diapositiva Final ---
\begin{frame}
    \centering
    \Huge \textbf{¿Preguntas?}
    
    \vspace{0.8cm}
    \large
    \textbf{Próxima clase:} Arquitectura Transformer: Visión General Encoder-Decoder
    
    \vspace{0.5cm}
    \normalsize
    Prof. Luis Alfredo Alvarado Rodríguez\\
    \small Escuela de Ciencias Físicas y Matemáticas
    
    \vspace{0.5cm}
    \footnotesize
    \textit{``Attention is all you need.''\\--- Vaswani et al., 2017}
\end{frame}

\end{document}